\section{Формализм}

\subsection{Определение ряда}
Обозначим уровни латинскими буквами \(a, b, \dots, z\) --- всего \nlevels{} уровней ---
и запишем значение \(n\)-го уровня как \(L(n)\), так что \(L(1)=a\) и \(L(\nlevels)=z\).
Ряд задаётся начальным значением и рекуррентным соотношением:
\[
  a = L(1) = \avalue, \qquad L(n+1) = L(n)^{\,L(n)}.
\]
Это (сдвинутая) \emph{тетрация}: основание степени само растёт с каждым шагом.

\subsection{Крипто кома}
\emph{Крипто комой} мы называем самовозведение последнего уровня:
\[
  \boxed{\;\ccoma = z^{z} = L(\nlevels)^{\,L(\nlevels)}.\;}
\]

\subsection{Компактная запись: оператор самовозведения}
Вся лестница --- это одна операция, применённая много раз. Введём \emph{оператор
самовозведения} \(\star\):
\[
  x^{\star} = x^{x} \qquad (\text{одна звёздочка} = \text{«возвести число в себя»}).
\]
Обозначая \(\star^{k}\) как \(k\)-кратное применение, всю лестницу и её вершину можно
записать одной строкой:
\[
  a = \avalue, \qquad L(n) = \bigl(\avalue\bigr)^{\star\,(n-1)}, \qquad
  \boxed{\;\ccoma = \bigl(\avalue\bigr)^{\star\,\nlevels}.\;}
\]
Иначе говоря: \(\avalue\), возведённое в себя \nlevels{} раз --- по одной звёздочке на
каждую букву алфавита (у \(a\) звёздочек нет, у \(z\) их 25, у крипто комы \(=z^{\star}\)
их \nlevels).

\subsection{Оценка величины (в логарифмической шкале)}
Уже \(b = a^{a} = \bigl(\avalue\bigr)^{\avalue} = \bvalue\); записать это число цифрами
невозможно. Поэтому мы работаем с \emph{порядком} \(M(n)=\log_{10}L(n)\). Логарифмируя
рекуррентное соотношение, получаем удобную формулу для порядка:
\[
  M(1) = \startlog, \qquad M(n+1) = M(n)\cdot 10^{\,M(n)}.
\]
Так, \(M(2) = \bmagnitude\), а начиная с уровня \(c\) сам порядок перестаёт помещаться в
машинное число и представляется \emph{башней} степеней десятки. Крипто кома \(\ccoma\)
есть степенная башня из десяток высотой примерно \ccheight{} этажей.

\subsection{Символ}
Для крипто комы мы предлагаем символ \(\ccoma\) --- рукописную «C» (Crypto Coma),
перечёркнутую двойной восходящей стрелкой \(\uparrow\uparrow\), отсылающей к восходящей
башне степеней (и к нотации Кнута для тетрации).

\subsection{«Мощность»}
В духе теории множеств мы оцениваем крипто кому как \emph{мощность} и предлагаем её
вниманию научного сообщества и всех, кто интересуется числами. Строго говоря, \(\ccoma\)
--- конечное натуральное число; о тонкости этого словоупотребления см.
раздел~\ref{sec:disclaimer}.
